% =============================================================================
% NOTATION TABLE
% =============================================================================
%
% PURPOSE:
%   Rigorously define every symbol appearing in a mathematical expression,
%   equation, or function.
%
% USAGE:
%   Place a notationtable immediately after the equation(s) it describes.
%   Include exactly ONE symbol per row.
%   Use template macros such as \reals, \SO, and \SE for mathematical types.
%   Vector and matrix shortcuts (e.g. \x, \A) require a shorthand environment.
%
%   \begin{notationtable}
%       $\theta$ & $\reals$ & Angle (rad) & Input angle. \\
%       $d$ & $\reals$ & Distance (m) & Output distance. \\
%       $f$ & $\reals \to \reals$ & Angle $\to$ Distance
%           & Maps an angle to a distance. \\
%   \end{notationtable}
%
% COLUMN DEFINITIONS:
%   1. Sym. (Symbol)
%      - Use the exact mathematical symbol appearing in the equation.
%      - ONE symbol per row; preserve subscripts, superscripts and formatting.
%   2. Type (Mathematical Type)
%      - Specify space, dimensions, or function signature.
%      - Examples: $\reals$, $\reals^n$,
%        $\reals^{m\times n}$, $\SO(3)$, $\SE(3)$,
%        $\reals^n \to \reals^m$.
%      - Do not substitute units or physical meaning for mathematical types.
%   3. Meaning (Semantic Type)
%      - State what the object represents, including units if applicable.
%      - For functions, explicitly give the semantic input/output mapping:
%        Angle -> Distance; Position -> Velocity; Image -> Descriptor.
%      - Distinguish meanings even when mathematical types are identical.
%      - Include coordinate frames when relevant.
%   4. Description
%      - Precisely define the object and its role in the equation.
%      - State reference frames, indexing conventions, and assumptions
%        as necessary. Explain what functions compute.
%
% GENERAL RULES FOR AUTHORS AND CODING AGENTS:
%   - ONE symbol per row; do not group distinct symbols.
%   - Define every symbol in the associated equation(s), including indices
%     and dimension parameters unless already defined elsewhere.
%   - Preserve notation consistently with the equations.
%   - Separate mathematical representation from semantic meaning.
%   - Include units for physical quantities whenever applicable.
%   - Specify coordinate frames for geometric quantities when relevant.
%   - Be concise but mathematically precise.
%   - Do not invent unknown dimensions, units, or interpretations.
%   - Put mathematical expressions in math mode.
%   - Do not add headers or horizontal rules; the environment supplies them.
%   - Tables can span multiple pages and are not floats.
%
% REQUIRED PACKAGES (load in the document preamble):
%   \usepackage{array}
%   \usepackage{longtable}
%   \usepackage{booktabs}
%   \usepackage{amsmath}
%   \usepackage{amssymb}
%
% =============================================================================

%TC:envir notationtable [] 0
\newlength{\notationtablewidth}
\newenvironment{notationtable}{%
    \small
    \renewcommand{\arraystretch}{1.15}%
    \setlength{\LTpre}{4pt}%
    \setlength{\LTpost}{6pt}%
    % Four p-columns with @{} at both ends have six tabcolsep paddings.
    \setlength{\notationtablewidth}{\dimexpr\linewidth-6\tabcolsep\relax}%
    \begin{longtable}{@{}
        >{\raggedright\arraybackslash}p{0.08\notationtablewidth}
        >{\raggedright\arraybackslash}p{0.20\notationtablewidth}
        >{\raggedright\arraybackslash}p{0.22\notationtablewidth}
        >{\raggedright\arraybackslash}p{0.50\notationtablewidth}@{}}
    \toprule
    \textbf{Sym.} & \textbf{Type} & \textbf{Meaning} & \textbf{Description} \\
    \midrule
    \endfirsthead
    \toprule
    \textbf{Sym.} & \textbf{Type} & \textbf{Meaning} & \textbf{Description} \\
    \midrule
    \endhead
    \midrule
    \multicolumn{4}{r@{}}{\footnotesize Continued on next page} \\
    \endfoot
    \bottomrule
    \endlastfoot
}{%
    \end{longtable}%
}

% Mathematical function signature, e.g. \maptype{\reals^n}{\reals^m}
\newcommand{\maptype}[2]{\mbox{\(#1 \to #2\)}}

% Semantic mapping, e.g. \semanticmap{Angle}{Distance}
\newcommand{\semanticmap}[2]{\mbox{#1 \(\to\) #2}}

% =============================================================================
% EXAMPLES (paste into the document body; do not use in the preamble)
% =============================================================================
%
% Point-coordinate transformation:
% \begin{shorthand}
% \[
%     \p_A = \R_{AB}\p_B + \t_{AB}.
% \]
% \begin{notationtable}
%     $\p_A$ & $\reals^3$ & Position (m)
%         & Coordinates of the point expressed in frame $A$. \\
%     $\p_B$ & $\reals^3$ & Position (m)
%         & Coordinates of the same point expressed in frame $B$. \\
%     $\R_{AB}$ & $\SO(3)$ & Rotation
%         & Rotation mapping vector coordinates from frame $B$ into frame $A$. \\
%     $\t_{AB}$ & $\reals^3$ & Translation (m)
%         & Position of the origin of frame $B$, expressed in frame $A$. \\
% \end{notationtable}
% \end{shorthand}
%
% Function example:
% \[
%     d = f(\theta).
% \]
% \begin{notationtable}
%     $f$ & \maptype{\reals}{\reals}
%         & \semanticmap{Angle}{Distance}
%         & Computes distance from angle according to a specified model. \\
%     $\theta$ & $\reals$ & Angle (rad)
%         & Input angle to $f$. \\
%     $d$ & $\reals$ & Distance (m)
%         & Distance computed by $f(\theta)$. \\
% \end{notationtable}
